\section{Introduction}\label{sec:intro}

ARC-AGI-3 is a benchmark of interactive games \citep{arcprize2026arcagi3}.
Each game is a grid of coloured cells with several levels (six to ten in the public games) and no instructions: the agent sees the board as an image, sends actions such as a move or a click, and has to work out the goal from what changes.
A level scores higher the fewer actions the agent needs compared with a human who played it, and later levels weigh more (Sec.~\ref{sec:setting}).
The competition plays 110 hidden games offline, on one 96 GB GPU, in nine hours; each team gets one scored run a day, which we call a \emph{draw}, and a draw returns a score and no log \citep{arcprize2026competitions}. % G7 A3 V3-W3

We ran the same program twice on the 25 public games, which the organisers call the \emph{demonstration set}, and it cleared 123 and then 101 levels. % A3 B1
Nothing in the code differed.
Almost every change we tried was smaller than that gap, so this paper is first about telling a change from noise, and then about what the instruments showed.

We asked whether, on one offline GPU with a fixed clock, the agent was limited by its decisions or by its tokens.
The answer so far is partial: serving gains were too small to see, the one level gain we found across days did not reproduce when its arms were interleaved, four changes to what the model is shown, when it must act and what it can check fired as designed on runs too few to resolve a level gain, and teaching the model from its own wins lowered levels.

The problem is hard for a plain reason: a run's score sums nearly two hundred all-or-nothing level outcomes, and one run per arm cannot tell a real gain from their spread.
Published ARC-AGI-3 agents use hosted models on the demonstration set \citep{lehmann2026tycho,skoutnev2026twin}, and the open-weight write-ups we read report no measurement protocol \citep{bessis2026duck,franzen2026writeup}.
One public repository, ``Before Tuning the Harness'', repeats identical submissions before tuning and accounts cost per request \citep{guoyi2026repeat}; we read its page and README only and did not check its model, and we price tokens in levels. % V3-W39
Our entry forks Daniel Franzen's public notebook \citep{franzen2026milestone2}; Sec.~\ref{sec:setting} separates what we inherited, borrowed and built.
We make five contributions.

\begin{enumerate}\setlength{\itemsep}{2pt}
\item \textbf{A measurement kit for a noisy fixed-budget agent (Sec.~\ref{sec:measure}).}
Six identical pairs give a run-to-run SD of 8.25 levels across days (5.6 around their shared first-to-second shift) and eight runs on one day 5.95, so two runs per arm resolve about 23 levels across days and 17 on one day by normal approximation; a $t$ test needs 5 to 12 runs per arm to see a 10-level gain, and a wave of twelve runs per arm made side by side would resolve about 9. % B1 B4 V3-X12 V3-W47 V3-W49
We also used concurrent replication on rented copies of the scoring GPU, which match Kaggle within 1\% on serving counters, and once a paired decision test, which replays single decisions with and without a change against the model's own sampling noise and is not yet validated against a lever with a known level effect. % V3-T3 V3-T4 V3-X1

\item \textbf{A token-budget curve that prices throughput (Sec.~\ref{sec:curve}).}
Built from recorded level attempts, it puts the agent at about 90\% of its ceiling, the score that unlimited tokens would give at the depths we recorded. % C2 N2
Our serving speed-ups are then worth about 4 levels (2 to 10 across our estimates), below what two runs per arm can see; the price is a simulator's and is itself uncertain. % C4 N1.3
Read from the clock, the last tenth of the window earns 7\% of the points, about three times the curve's price on score. % V3-W13 V3-W14

\item \textbf{A context lever that failed an interleaved test (Sec.~\ref{sec:yield}).}
Yielding a turn after 8,192 generated tokens, four times Franzen's threshold, re-sends the turn's opening message and board image about 70\% less often. % D14 D5
Across days it coincided with 12.5 more levels. % D1 D8
Interleaved on builds that differ in one line, six runs per arm, it scored 2.5 fewer ($p$ 0.42), which excludes the earlier gain but not small gains: drift between days and extra tokens are the candidate causes, and only the interleaved design could show that the gain did not reproduce (Sec.~\ref{sec:yield}). % YN6 V3-W36 V3-W38

\item \textbf{A diagnosis-first harness program with kill rules (Sec.~\ref{sec:harness}).}
In eight paired plays, stuck plays often stopped acting, wrote off object types untested and misread what the harness shows; over 562 level plays of the same four runs, 27\% of never-cleared levels hold three or more idle turns in a row, against 8\% of cleared ones. % V3-W18
Five changes, each gated by a rule in a dated file written before its runs, fired as designed: a stall check, a corner-joined view of objects, a scroll detector, tool fixes that keep a timed-out snippet's output, and a replay check that tests the model's simulator and goal guesses against the level's history at no cost in actions. % V3-W20 V3-W22 V3-W23 V3-W28 V3-W29 V3-W38
The first two led a pooled control by about five levels and the replay check by about three (over twice as much against the two Kaggle controls alone), which three runs cannot resolve; the scroll detector gave no gain and was not shipped; the replay check's helper reported a mismatch in 52 of 85 replays (42 a wrong prediction), which counts reports and not better play. % V3-W19 V3-W23 V3-W29 V3-W41
None adds tokens, so this is evidence about decisions, not budget.
A research pass reads why such changes are the place to look: on long levels most of the clock goes to finding the rule or the goal, not to carrying out a plan. % V3-W31 V3-W34

\item \textbf{Two fine-tuning results and what came with them (Sec.~\ref{sec:finetune}).}
(a) Rejection-sampling fine-tuning on the agent's own cleared levels lowered validation loss by 13.9\% and lowered levels. % V3-L6 V3-L9 V3-L10
Its three runs sit at or below the two lowest of 18 runs without it (rank $p$ 0.004, two adapter scales pooled). Two later same-card controls without it cleared more, and the adapter's three runs are the lowest of those five (rank $p$ 0.10, the smallest possible for three against two). The adapter and the controls ran on different services and days, not interleaved, and the controls are two, one restarted after a preemption. % V3-W9 V3-W44
The agent already clears 91\% of the training games' levels, so the filter kept nearly everything, slow wins included; the adapter's reasoning per action fell toward the training games' level, and its resets and deaths rose. % V3-L1 V3-L2 V3-L13 V3-L18
This chain is our best reading of why; two alternatives are untested.

(b) In reinforcement learning on this stack the training forward pass is not deterministic: two identical passes agree exactly on about 12\% of tokens and differ by 0.37 nats per token (RMS). % V3-G5
A drift guard must therefore be read net of the stack's own floor; read raw, our predeclared guard would have discarded the run on noise. % V3-G4 V3-G7
Round A shipped no adapter: a first-wave lead of 1.9 levels, selection-biased since only the leader went on, did not replicate (second wave $-2.8$, pooled $+0.25$ against a $+1.75$ bar). A round B adapter was nonetheless pinned for three board draws by our override before its gate read (35.90, 39.79, 35.98) and then failed (Sec.~\ref{sec:rl}). % V3-W45 V3-W17
\end{enumerate}

\begin{figure}[t]
\centering
\figfile{fig_loop}{3.4in}
\caption{One turn of the agent, and where each lever acts. A turn ends when the model acts. A yield ends a request after a set number of tokens generated without an action and resumes the turn with the opener (the turn's opening message) and the board image sent again. Accent: the three key decisions of Sec.~\ref{sec:setting} (the 8,192-token yield; expert pruning that frees memory for 14 concurrent games; the fail-closed checks, which have not yet fired on a scored run); gray: inherited parts. Schematic.}\label{fig:loop} % A4 D5
\end{figure}

\paragraph{What we do not claim.}
We attribute no board draw to any lever, and no level gain to the yield or to the harness changes.
Numbers from the demonstration set are development numbers, not held out.
Our 40.98 is the best of 28 live draws; the agent's seventeen valid draws average 34.59 (32.67 with the 0.00 of unknown cause). The best of 17 to 21 such draws is expected near 40 to 41, so 40.98 is no evidence for one lever, and we expect a private score of about 35 to 37, not 40.98 (Appendix~\ref{app:noisegates}). % V3-W16 V3-W30 V3-W52
We make no claim about whether the model's licence is eligible for a prize. % H7

% count:off  (the Accuracy line is a data line, not prose)
\notebox{\textbf{Result.} Public score \textbf{40.98} on half of the 110 hidden games: Kaggle submission \subid{57417773}, drawn 2026-10-20\pubonly{, \rankof{26}{5,073} teams on the final public board at the 2026-11-02 deadline}. It is the best of 28 live draws since the fork; the agent's seventeen valid draws (errors, zeros, adapter draws and field-wide failure days excluded) span 28.97 to 40.98 (mean 34.59; 32.67 with the 0.00 of 10-18). The private score is pending.} % V3-B1 V3-B2 V3-B9 V3-W30 V3-W43 BOARD-REFRESH: once posted, lead this line with the private score of the linked Final Submission.
% count:on
